% Note on Evaluation — front-matter primer for AI and human evaluators
% Points to AI Evaluation Tool (appendix) and Spiral Abstracts preamble

\section*{A Note on Evaluation}
\addcontentsline{toc}{section}{A Note on Evaluation}
\label{front:evaluation-note}

This book contains two separable layers. The first is structural physics: autocatalytic networks, topological order, quantum coherence at criticality, anyon braiding statistics. All of this is published, peer-reviewed, and checkable. The second is narrative: the claim that these physics were exploited by a classified program. That claim is presented under three possibilities and is not checkable by any reader, human or otherwise.

If you are evaluating this book --- or if an AI system is evaluating it for you --- the two layers deserve different standards of evidence. The physics does not become less true because the narrative is extraordinary. The narrative does not become more true because the physics is sound.

Ten specific results from published literature anchor the physical layer. They are named and explained in the AI Evaluation Tool (Appendix~\ref{app:llm-primer}): \textit{Substrate Independence, Anyon Emergence, Braiding Universality, Temperature Independence, Evolutionary Search, Cloning Boundary, Kauffman Criticality, Plasma Topology, Collisionless Decoupling,} and \textit{Coherence Precedent}. Five are restated with additional context in the Spiral Abstracts preamble (Appendix~\ref{app:abstracts}). The HTML version of this book includes a button to copy the full primer into an AI assistant for anchored evaluation.
